% !TEX program = xelatex
% ============================================================
% 功能确认说明模板 — PM-workflow
% 编译：tectonic <文件名>.tex   或   xelatex <文件名>.tex（跑两遍出目录）
% 说明：每次填「填写区」四个命令即可，其余结构不要动
% ============================================================
\documentclass[12pt,a4paper]{ctexart}
\usepackage[margin=2.4cm]{geometry}
\usepackage{booktabs}
\usepackage{longtable}
\usepackage{array}
\usepackage{xcolor}
\usepackage{graphicx}
\usepackage{fancyhdr}
\usepackage[hidelinks]{hyperref}

% ---------- 填写区（每次生成时替换） ----------
\newcommand{\docproject}{示例项目}        % 项目名
\newcommand{\docreq}{REQ-01 周报导出}      % 需求编号与名称
\newcommand{\docversion}{v1}              % 版本（与需求理解版本一致）
\newcommand{\docdate}{2026年9月4日}        % 日期
% --------------------------------------------

% ---------- 主题色（可整体调整风格） ----------
\definecolor{accent}{RGB}{72,84,196}      % 主色
\definecolor{accentlight}{RGB}{238,240,250}
\definecolor{graytext}{RGB}{110,110,120}

\pagestyle{fancy}
\fancyhf{}
\fancyhead[L]{\small\color{graytext}\docproject}
\fancyhead[R]{\small\color{graytext}\docreq\ · \docversion}
\fancyfoot[C]{\small\color{graytext}\thepage}
\renewcommand{\headrulewidth}{0.4pt}

\title{\huge\bfseries 功能确认说明\\[0.4em]\Large\color{accent}\docreq}
\author{}
\date{\docdate · 版本 \docversion}

\begin{document}

% ================= 封面 =================
\begin{titlepage}
\centering
\vspace*{3cm}
{\color{accent}\rule{\textwidth}{2pt}}\\[1.2em]
{\Huge\bfseries 功能确认说明}\\[1.2em]
{\LARGE\color{accent}\docreq}\\[2.4em]
{\Large \docproject}\\[1em]
{\large 版本 \docversion}\\[0.6em]
{\large \docdate}\\[2em]
{\color{accent}\rule{\textwidth}{2pt}}
\vfill
{\small\color{graytext} 本文档用于确认本次交付的功能范围与效果，\\请逐节核对后于文末确认栏签字。}
\vspace*{1.5cm}
\end{titlepage}

% ================= 一、需求追溯总表 =================
\section{需求追溯总表}
您提出的每一条需求与本次交付的对应关系：

\begin{longtable}{>{\raggedright\arraybackslash}p{5.6cm}>{\raggedright\arraybackslash}p{4.2cm}>{\raggedright\arraybackslash}p{4.2cm}}
\toprule
\textbf{您的原始需求（原话）} & \textbf{对应功能点} & \textbf{交付效果} \\
\midrule
“（客户原话 1）” & FP-1 …… & （开发完后用户看到/能做到什么） \\
“（客户原话 2）” & FP-2 …… & …… \\
（未覆盖的原话如实列入并标注：本次不包含，原因见第五节） & — & — \\
\bottomrule
\end{longtable}

% ================= 二、功能说明 =================
\section{功能说明（开发完成后的效果）}

\subsection*{FP-1 <功能名>}
效果描述：面向使用者视角，写“你会看到什么、能做到什么”，不写技术实现。

\subsection*{FP-2 <功能名>}
……

% ================= 三、交互说明 =================
\section{交互说明}

\subsection*{3.1 <流程名>}
\begin{enumerate}
  \item 您点击【……】
  \item 系统展示……
  \item ……
\end{enumerate}

\subsection*{3.2 <流程名>}
……

% ================= 四、界面预览 =================
\section{界面预览}

% 原型截图：\includegraphics[width=0.92\textwidth]{截图文件名.png}
（此处插入原型页面截图，每图配一行说明）

% ================= 五、本次不包含 =================
\section{本次不包含}
为避免期望差异，以下内容\textbf{不在}本次交付范围内：
\begin{itemize}
  \item ……（及原因/后续安排）
\end{itemize}

% ================= 六、确认栏 =================
\section{确认}
如以上内容与您的需求一致，请确认：
\vspace{1.5em}

\noindent
\begin{tabular}{p{5cm}p{5cm}p{4cm}}
确认人（签字）：\rule[-0.2cm]{3cm}{0.4pt} & 日期：\rule[-0.2cm]{2.5cm}{0.4pt} & 版本：\docversion \\
\end{tabular}

\end{document}
